\documentclass[11pt]{article}

\usepackage[T1]{fontenc}
\usepackage{lmodern}
\usepackage[margin=1.03in]{geometry}
\usepackage{amsmath,amssymb,amsthm,mathtools}
\usepackage{microtype}
\usepackage{graphicx}
\usepackage[hidelinks]{hyperref}
\usepackage{enumitem}

\setlist{nosep}
\newtheorem{theorem}{Theorem}
\newtheorem{lemma}{Lemma}
\theoremstyle{remark}
\newtheorem{remark}{Remark}

\newcommand{\Ba}{\operatorname{Ba}}
\newcommand{\R}{\mathbb{R}}
\newcommand{\Q}{\mathbb{Q}}
\newcommand{\N}{\mathbb{N}}
\newcommand{\norm}[1]{\lVert #1\rVert}

\title{A Dense-Set Counterexample for Weak-Baire Measurability of the Norm}
\author{Independent research note}
\date{August 13, 2026}

\begin{document}
\maketitle

\begin{abstract}
Question~(B) of Avil\'es--Plebanek--Rodr\'iguez asks whether relative
weak-Baire measurability of a Banach-space norm on a norm-dense set forces
weak-Baire measurability on the whole space.  The answer is no in ZFC.
For every nonseparable real Hilbert space \(X\), we construct a norm-dense
subset \(A\subseteq X\) such that the restricted norm is the trace of a
weak-Baire function depending on one linear coordinate, although the norm on
\(X\) is not weak-Baire measurable.  The parenthetical variant in which
\(A\) must be a dense linear subspace is not resolved here.
\end{abstract}

\section*{Packet status}

\begin{description}[leftmargin=3.0cm,style=nextline]
\item[Run] \texttt{fa\_banach\_001}, lane 13.
\item[Source] A.~Avil\'es, G.~Plebanek, and J.~Rodr\'iguez,
\emph{A weak\(^{*}\) separable \(C(K)^{*}\) space whose unit ball is not
weak\(^{*}\) separable}, arXiv:1112.5710, Section~3.3, Question~(B),
PDF page~22 \cite{APR}.
\item[Classification] Candidate counterexample, likely valid subject to
human mathematical review.
\item[Scope] Full negative answer to the norm-dense-set formulation; no
claim for the dense-linear-subspace parenthetical.
\end{description}

The source question is reproduced here:

\begin{center}
\includegraphics[width=0.96\textwidth]{figures/open_problem_crop.png}
\end{center}

In transcription: if \(A\subseteq X\) is norm dense and the norm is
relatively \(\Ba(X,w)\)-measurable on \(A\), must the norm be
\(\Ba(X,w)\)-measurable on \(X\)?  The paper also asks the stronger-looking
parenthetical variant where \(A\) is a linear subspace.

\section{Proof intuition}

Two elementary ideas fit together.  First, every real function measurable
for \(\Ba(X,w)=\sigma(X^*)\) depends on only countably many dual
coordinates.  Countably many vectors cannot norm a nonseparable Hilbert
space: a nonzero vector orthogonal to all of them is indistinguishable from
zero.  Hence the global Hilbert norm is not weak-Baire measurable.

Second, a dense subset need not retain any linear geometry.  Split a Hilbert
space as \(\R\oplus_2 Z\).  Choose a Borel function
\(\psi:\R\to[0,\infty)\) with dense graph and retain from each vertical
fiber only the sphere \(\norm{z}=\psi(t)\).  The resulting set is dense, while
its norm is \(\sqrt{t^2+\psi(t)^2}\), a Borel function of the one weak
coordinate \(t\).  Thus the restriction is measurable even though the
global norm is not.

\section{The counterexample}

We first record the countable-dependence fact in the needed form.

\begin{lemma}[Countable-coordinate dependence]\label{lem:countable}
Let \(\mathcal F\) be any family of maps from a set \(\Omega\) to \(\R\).
If \(u:\Omega\to\R\) is measurable for \(\sigma(\mathcal F)\), then there is
a countable subfamily \(\mathcal F_0\subseteq\mathcal F\) such that
\[
 f(x)=f(y)\quad(f\in\mathcal F_0)
 \qquad\Longrightarrow\qquad u(x)=u(y).
\]
\end{lemma}

\begin{proof}
Call a subset of \(\Omega\) countably determined if it belongs to
\(\sigma(\mathcal F_0)\) for some countable
\(\mathcal F_0\subseteq\mathcal F\).  The countably determined sets form a
sigma-algebra: for a countable union, take the union of the countably many
subfamilies involved.  This sigma-algebra contains every generator, so it is
all of \(\sigma(\mathcal F)\).

For each \(r\in\Q\), choose a countable subfamily determining
\(u^{-1}((-\infty,r))\), and take their union \(\mathcal F_0\).  Points having
the same \(\mathcal F_0\)-coordinates belong to exactly the same rational
sublevel sets of \(u\), and therefore have the same \(u\)-value.
\end{proof}

\begin{theorem}\label{thm:main}
Let \(X\) be any nonseparable real Hilbert space.  There exists a norm-dense
set \(A\subseteq X\) such that \(\norm{\cdot}|_A\) is measurable for the
trace of \(\Ba(X,w)\) on \(A\), while \(\norm{\cdot}:X\to\R\) is not
\(\Ba(X,w)\)-measurable.
\end{theorem}

\begin{proof}
Choose a one-dimensional summand and write
\[
 X=\R\oplus_2 Z,
 \qquad \norm{(t,z)}=\sqrt{t^2+\norm{z}^2},
\]
where \(Z\) is nonseparable.  Let \((q_n)_{n\ge1}\) enumerate
\(\Q\cap[0,\infty)\), put
\[
 D_n=n\sqrt2+\Q\quad(n\ge1),
 \qquad
 D_0=\R\setminus\bigcup_{n\ge1}D_n,
\]
and define
\[
 \psi(t)=
 \begin{cases}
 q_n,&t\in D_n,\ n\ge1,\\
 0,&t\in D_0.
 \end{cases}
\]
The sets \(D_n\) are pairwise disjoint, countable, dense, and Borel; their
complement \(D_0\) is also Borel.  Thus \(\psi\) is Borel.  Its graph is
dense in \(\R\times[0,\infty)\): every open rectangle contains
\((t,q_n)\) after first choosing a rational \(q_n\) in its vertical interval
and then \(t\in D_n\) in its horizontal interval.

Define
\[
 A=\{(t,z)\in\R\oplus_2 Z:\norm{z}=\psi(t)\}.
\]
To see that \(A\) is norm dense, fix \((t,z)\in X\) and
\(\varepsilon>0\).  Choose \(q_n\) with
\(\lvert q_n-\norm z\rvert<\varepsilon/2\), and choose
\(s\in D_n\) with \(\lvert s-t\rvert<\varepsilon/2\).  If \(z\ne0\), set
\(z'=q_nz/\norm z\); if \(z=0\), choose \(z'\in Z\) with
\(\norm{z'}=q_n\).  Then \((s,z')\in A\) and
\[
 \norm{(s,z')-(t,z)}
 =\sqrt{\lvert s-t\rvert^2+\lvert q_n-\norm z\rvert^2}<\varepsilon.
\]

Let \(\pi:X\to\R\) be the first-coordinate functional and put
\(\phi(t)=\sqrt{t^2+\psi(t)^2}\).  The function \(\phi\) is Borel, hence
\(\phi\circ\pi\) is \(\Ba(X,w)\)-measurable.  For every \((t,z)\in A\),
\[
 (\phi\circ\pi)(t,z)=\sqrt{t^2+\psi(t)^2}
 =\sqrt{t^2+\norm z^2}=\norm{(t,z)}.
\]
Therefore the restricted norm on \(A\) is measurable for the trace of
\(\Ba(X,w)\).

It remains to show that the global norm is not weak-Baire measurable.  For a
Banach space, \(\Ba(X,w)=\sigma(X^*)\).  If the norm were measurable,
Lemma~\ref{lem:countable} would give \(x_n^*\in X^*\), \(n\ge1\), on whose
joint coordinates it depends.  By the Riesz representation theorem, the
\(x_n^*\)'s are represented by a countable family of Hilbert vectors.  Their
closed linear span is separable and hence proper in nonseparable \(X\).
Choose \(0\ne v\) orthogonal to it.  Then \(x_n^*(v)=x_n^*(0)\) for every
\(n\), so countable dependence gives \(\norm v=\norm0=0\), a contradiction.
\end{proof}

\section{Scope, verification, and novelty}

The construction is ZFC and already works in
\(X=\R\oplus_2\ell_2(\omega_1)\).  It answers the source implication
negatively when \(A\) is an arbitrary norm-dense set.  It deliberately does
not make \(A\) linear.  Requiring linearity turns the construction into the
dense graph of a discontinuous operator over a countably coordinatized
quotient; no valid construction or impossibility proof for that variant was
found in the upgrade audit.

The proof was checked at four pressure points: the cosets \(n\sqrt2+\Q\) are
pairwise disjoint; the graph of \(\psi\) is dense; radial rescaling gives the
displayed exact distance formula; and Lemma~\ref{lem:countable} applies to
the rational sublevel sets of the norm.  No computational assumption or
set-theoretic hypothesis beyond ZFC is used.

A bounded search on August 13, 2026 covered the run indexes, arXiv:1112.5710,
the exact text of Question~(B), the phrases \emph{relatively
\(\Ba(X,w)\)-measurable}, \emph{norm dense set}, and \emph{weak Baire norm},
plus the title and authors.  It found the source and nearby work on Baire
structures in \(C(K)\), but no later explicit answer and no matching shell
construction.  Novelty confidence is moderate because the proof is
elementary and may be folklore.  Human review should focus on literature
novelty and on the intended parsing of the source's parenthetical
\emph{(or linear subspace)}; the mathematical counterexample itself is
self-contained.

\begin{thebibliography}{9}
\bibitem{APR}
A.~Avil\'es, G.~Plebanek, and J.~Rodr\'iguez,
\emph{A weak\(^{*}\) separable \(C(K)^{*}\) space whose unit ball is not
weak\(^{*}\) separable}, arXiv:1112.5710 (2011).
\end{thebibliography}

\end{document}

